\documentclass[11pt]{article}
\usepackage[margin=1in]{geometry}
\usepackage{amsmath,amssymb,amsthm}
\usepackage{graphicx}
\usepackage{booktabs}
\usepackage{xcolor}
\usepackage{listings}
\usepackage[numbers]{natbib}
\usepackage{hyperref}

\newtheorem{theorem}{Theorem}
\newtheorem{lemma}[theorem]{Lemma}
\theoremstyle{definition}
\newtheorem{definition}[theorem]{Definition}
\theoremstyle{remark}
\newtheorem{remark}[theorem]{Remark}

\lstdefinelanguage{lean}{
  keywords={theorem,lemma,def,abbrev,namespace,end,open,import,by,intro,
    rintro,obtain,have,set,with,fun,at,if,then,else,show,exact,refine,apply,
    rw,simp,simpa,unfold,cases,rcases,by_cases,omega,ring,ext,left,right,constructor},
  sensitive=true,
  comment=[l]{--},
  morecomment=[s]{/-}{-/},
  string=[b]"
}
\lstset{language=lean,basicstyle=\ttfamily\footnotesize,breaklines=true,
  columns=fullflexible,keepspaces=true,
  keywordstyle=\color{blue!60!black}\bfseries,commentstyle=\color{green!40!black},
  frame=single,
  literate=
    {·}{{$\cdot$}}1 {×}{{$\times$}}1 {₁}{{$_1$}}1 {₂}{{$_2$}}1
    {ℕ}{{$\mathbb{N}$}}1 {ℤ}{{$\mathbb{Z}$}}1 {←}{{$\leftarrow$}}1
    {→}{{$\to$}}1 {↔}{{$\leftrightarrow$}}1 {∀}{{$\forall$}}1
    {∃}{{$\exists$}}1 {∈}{{$\in$}}1 {∑}{{$\sum$}}1 {∨}{{$\lor$}}1
    {⊢}{{$\vdash$}}1 {⟨}{{$\langle$}}1 {⟩}{{$\rangle$}}1
}

\title{The 8-connected bit-quad Euler formula for finite rasters:\\
a Lean proof by vertex double counting}
\author{Justin Hart\\Viridis LLC}
\date{28 October 2026}

\begin{document}
\maketitle

\begin{abstract}
For a finite set $B\subset\mathbb{Z}^2$ of pixels, let $K(B)$ be the union of the
closed unit squares of $B$, regarded as a cubical complex with Euler
characteristic $\chi(K(B)) = V-E+F$. Gray's bit-quad formula expresses this
topological invariant through purely local counts:
$4\chi(K(B)) = Q_1 - Q_3 - 2Q_D$, where $Q_1$, $Q_3$ and $Q_D$ count the
$2\times2$ windows holding exactly one pixel, exactly three pixels, and exactly
two diagonally opposite pixels of $B$. We give a Lean~4 proof, against Mathlib,
of the sealed statement \texttt{P1.BitQuad}, which states this identity for every
finite raster. The proof is a double count over the vertices of $K(B)$. Each of
$F$, $E$ and $V$ is written as a sum over vertices of indicator functions of
translated sets. A single shift-counting lemma evaluates all eight such sums,
and a pointwise identity is closed by a case split on the 16 window patterns, of
which only the 15 non-empty ones can occur at a vertex. This is the first rung (L1) of a lemma ladder
whose goal is to replace the topological hypothesis assumed in the
machine-checked core of \emph{Trees Cannot Cut} with a theorem about rasters.
\end{abstract}

\section{Introduction}

Hart's \emph{Trees Cannot Cut}~\cite{Hart2026} argues that fuel breaks and
habitat corridors split a forest landscape into separate compartments only
through their cycles: an acyclic corridor leaves the fuel matrix connected. The
graph theory behind this claim is machine-checked, but the step from topology to
graph theory is carried as an explicit hypothesis (\texttt{FaceCount}), which
says that the number of compartments equals the circuit rank plus one. The
manuscript justifies it by Alexander duality. Mathlib~\cite{Mathlib2020} has
neither Alexander duality nor a planar Euler formula in a form that applies
here.

Fire and fuel-break data usually come as rasters, so project P1 aims to
discharge the hypothesis in the raster setting. For a finite barrier raster $B$
with 8-connected barriers and a 4-connected complement, the target identity is
\[
  \#\{\text{enclosed 4-components of } \mathbb{Z}^2\setminus B\}
  \;=\; C_8(B) - \chi(K(B)).
\]
The proposed route is a ladder of three lemmas. L1 is a local formula for
$\chi(K(B))$. L2 says that adding one pixel preserves
$\text{enclosed} - C_8 + \chi$. L3 is the bridge itself, by induction on $|B|$.
This paper does L1: it gives a Lean proof of Gray's bit-quad formula in its
8-connected form. L1 matters for the rest of the ladder because it makes the
change in $\chi$ under a single-pixel insertion a purely local quantity: only
the four windows around the new pixel change.

\section{Background and related work}

Gray~\cite{Gray1971} introduced locally countable properties of binary images,
among them area, perimeter and the Euler number. He showed that the Euler number
can be computed from counts of $2\times2$ ``bit-quad'' patterns, with different
weightings for 4- and 8-connectivity of the foreground. The (8,4) pairing of
connectivity for foreground and background, and the role it plays in digital
Jordan-type theorems, are surveyed by Kong and Rosenfeld~\cite{KongRosenfeld1989}.
In the 8-connected reading, two pixels that share only a corner touch in $K(B)$.
The two diagonal patterns then contribute $-2$ to $4\chi$, because their
shared corner is one vertex where four edges meet. Our Lean development uses
Lean~4~\cite{Lean4} and Mathlib~\cite{Mathlib2020}. We are not aware of an
earlier machine-checked proof of the bit-quad formula, but we have not searched
exhaustively and make no claim of novelty.

\section{Setting: the sealed statement}

The statement module \texttt{ViridisStatements.P1.BitQuad} fixes the following
definitions. A pixel $p=(i,j)\in\mathbb{Z}^2$ stands for the closed square
$[i,i+1]\times[j,j+1]$.
\begin{itemize}
\item $\mathrm{vertices}(B)=\bigcup_{p\in B}\mathrm{corners}(p)$, where
  $\mathrm{corners}(i,j)=\{(i,j),(i+1,j),(i,j+1),(i+1,j+1)\}$.
\item $\mathrm{hEdges}(B)=\bigcup_{(i,j)\in B}\{(i,j),(i,j+1)\}$ (horizontal
  edges named by their left end) and
  $\mathrm{vEdges}(B)=\bigcup_{(i,j)\in B}\{(i,j),(i+1,j)\}$ (vertical edges named
  by their lower end).
\item The Euler characteristic, computed in $\mathbb{Z}$:
  \[\chi(B) = |\mathrm{vertices}(B)| - (|\mathrm{hEdges}(B)|+|\mathrm{vEdges}(B)|) + |B|.\]
\item The window at a lattice point $v=(x,y)$ consists of four bits,
  $a=[(x-1,y-1)\in B]$, $b=[(x,y-1)\in B]$, $c=[(x-1,y)\in B]$ and
  $d=[(x,y)\in B]$, with $n(v)=a+b+c+d$.
\item $Q_1$, $Q_3$ and $Q_D$ count the vertices $v$ with $n(v)=1$, with
  $n(v)=3$, and with $\{a,d\}$ or $\{b,c\}$ as the only set bits.
\end{itemize}
Every window that meets $B$ is centred on a vertex of $K(B)$, so counting over
the vertices counts every such window exactly once.

\begin{theorem}[\texttt{P1.BitQuad}]\label{thm:main}
For every finite $B\subset\mathbb{Z}^2$,
$4\,\chi(K(B)) = Q_1(B) - Q_3(B) - 2\,Q_D(B)$.
\end{theorem}

\section{Proof architecture}

Write $\mathcal V=\mathrm{vertices}(B)$, and use $[P]\in\{0,1\}$ for the indicator
of a proposition $P$.

\begin{lemma}[Edge membership]\label{lem:edges}
$(i,j)\in\mathrm{hEdges}(B) \iff (i,j)\in B \lor (i,j-1)\in B$, and
$(i,j)\in\mathrm{vEdges}(B) \iff (i,j)\in B \lor (i-1,j)\in B$.
\end{lemma}

\begin{lemma}[Shift counting]\label{lem:shift}
Let $S,\mathcal V\subset\mathbb{Z}^2$ be finite and $s\in\mathbb{Z}^2$. If $x+s\in\mathcal V$
for every $x\in S$, then $\sum_{v\in\mathcal V}[v-s\in S]=|S|$.
\end{lemma}
\begin{proof}
The map $v\mapsto v-s$ is a bijection from $\{v\in\mathcal V: v-s\in S\}$ onto
$S$. Its inverse is $x\mapsto x+s$, which lands in $\mathcal V$ by hypothesis.
\end{proof}

Lemma~\ref{lem:shift} does all of the double counting. Every pixel $x\in B$
has its four corners $x$, $x+(1,0)$, $x+(0,1)$ and $x+(1,1)$ in $\mathcal V$.
Applying the lemma with the four shifts gives
\[
  4F = \sum_{v\in\mathcal V} n(v).
\]
By Lemma~\ref{lem:edges}, both ends $x$ and $x+(1,0)$ of a horizontal edge are
corners of the pixel above or below it, and likewise for vertical edges.
Write the arms of $v$ as left $[v-(1,0)\in\mathrm{hEdges}]$, right
$[v\in\mathrm{hEdges}]$, down $[v-(0,1)\in\mathrm{vEdges}]$ and up
$[v\in\mathrm{vEdges}]$. Applying the lemma once more for each arm gives
\[
  2E = \sum_{v\in\mathcal V}\mathrm{arms}(v), \qquad V=\sum_{v\in\mathcal V}1 .
\]
Hence $4\chi = \sum_{v\in\mathcal V}\bigl(4-2\,\mathrm{arms}(v)+n(v)\bigr)$.
Lemma~\ref{lem:edges} expresses each arm through the window bits. The left arm
is $a\lor c$, the right arm $b\lor d$, the down arm $a\lor b$ and the up arm
$c\lor d$. The summand therefore depends only on $(a,b,c,d)$, and
Table~\ref{tab:patterns} evaluates it. Every vertex is a corner of some pixel,
which must be one of the four window cells, so the empty pattern never occurs.
On the other 15 patterns the summand equals
$[n=1]-[n=3]-2[\text{diagonal}]$. Summing over $\mathcal V$ proves
Theorem~\ref{thm:main}.

\begin{table}[h]
\centering
\begin{tabular}{lcccc}
\toprule
Pattern class & \#patterns & $n$ & arms & $4-2\,\mathrm{arms}+n$\\
\midrule
empty (cannot occur) & 1 & 0 & 0 & 4\\
one pixel & 4 & 1 & 2 & $1$\\
adjacent pair & 4 & 2 & 3 & $0$\\
diagonal pair & 2 & 2 & 4 & $-2$\\
three pixels & 4 & 3 & 4 & $-1$\\
full & 1 & 4 & 4 & $0$\\
\bottomrule
\end{tabular}
\caption{The local summand on the 16 window patterns.}
\label{tab:patterns}
\end{table}

\section{The Lean development}

The candidate file \texttt{proofs/P1.BitQuad.lean} (145 lines) follows the
architecture above.
\begin{itemize}
\item \texttt{mem\_vertices}, \texttt{mem\_corners}, \texttt{mem\_hEdges} and
  \texttt{mem\_vEdges} unfold the membership predicates (Lemma~\ref{lem:edges}).
\item \texttt{card\_shift} proves Lemma~\ref{lem:shift} with
  \texttt{Finset.card\_bij}. Its corollary \texttt{sum\_shift} restates it as a
  sum of $\{0,1\}$-valued integer indicators via \texttt{Finset.sum\_boole}.
\item In \texttt{main}, the eight instances of \texttt{sum\_shift} (four for
  faces, two for each edge direction) and the three expressions of $Q_1$, $Q_3$
  and $Q_D$ as indicator sums reduce the goal to a single
  \texttt{Finset.sum\_congr}. The pointwise goal is closed by
  \texttt{by\_cases} on the four window memberships followed by \texttt{simp},
  which also uses the fact that the window at a vertex is non-empty.
\end{itemize}
The brief suggested evaluating the 16 patterns with \texttt{decide}. That
proved unnecessary: once the window memberships are case-split, \texttt{simp}
closes all 16 branches directly. The file uses no \texttt{decide} on large
objects and no \texttt{native\_decide}. The project checker
(\texttt{viridis-lean-check}) reports that the file compiles under Lean
\texttt{v4.28.0} with the pinned Mathlib, matches the sealed statement, depends
only on \texttt{propext}, \texttt{Classical.choice} and \texttt{Quot.sound}, and
replays through the kernel. Whether the claim is accepted is decided by the
project's independent verifier; see its certification page. The complete file
follows.

\begin{lstlisting}
import Mathlib
import ViridisStatements.P1.BitQuad

namespace Candidate

open ViridisStatements.P1.BitQuad

/-- Membership in `vertices`. -/
theorem mem_vertices (B : Finset Pixel) (v : ℤ × ℤ) :
    v ∈ vertices B ↔ ∃ p ∈ B, v ∈ corners p := by
  unfold vertices
  simp only [Finset.mem_biUnion]

theorem mem_corners (p v : ℤ × ℤ) :
    v ∈ corners p ↔ v = (p.1, p.2) ∨ v = (p.1 + 1, p.2) ∨ v = (p.1, p.2 + 1) ∨
      v = (p.1 + 1, p.2 + 1) := by
  unfold corners
  simp only [Finset.mem_insert, Finset.mem_singleton]

/-- A horizontal edge is present iff the pixel above or below it is in `B`. -/
theorem mem_hEdges (B : Finset Pixel) (h : ℤ × ℤ) :
    h ∈ hEdges B ↔ (h.1, h.2) ∈ B ∨ (h.1, h.2 - 1) ∈ B := by
  unfold hEdges
  simp only [Finset.mem_biUnion, Finset.mem_insert, Finset.mem_singleton]
  constructor
  · rintro ⟨p, hp, rfl | rfl⟩
    · left; simpa using hp
    · right; simpa using hp
  · rintro (h1 | h1)
    · exact ⟨_, h1, Or.inl rfl⟩
    · exact ⟨_, h1, Or.inr (by simp)⟩

/-- A vertical edge is present iff the pixel right or left of it is in `B`. -/
theorem mem_vEdges (B : Finset Pixel) (e : ℤ × ℤ) :
    e ∈ vEdges B ↔ (e.1, e.2) ∈ B ∨ (e.1 - 1, e.2) ∈ B := by
  unfold vEdges
  simp only [Finset.mem_biUnion, Finset.mem_insert, Finset.mem_singleton]
  constructor
  · rintro ⟨p, hp, rfl | rfl⟩
    · left; simpa using hp
    · right; simpa using hp
  · rintro (h1 | h1)
    · exact ⟨_, h1, Or.inl rfl⟩
    · exact ⟨_, h1, Or.inr (by simp)⟩

/-- Shift-counting: if every translate `x + (a, b)` of `S` lies in `V`, then
counting the `v ∈ V` with `v - (a, b) ∈ S` gives `|S|`. -/
theorem card_shift (V S : Finset (ℤ × ℤ)) (a b : ℤ)
    (hS : ∀ x ∈ S, (x.1 + a, x.2 + b) ∈ V) :
    (V.filter fun v => (v.1 - a, v.2 - b) ∈ S).card = S.card := by
  apply Finset.card_bij (fun v _ => (v.1 - a, v.2 - b))
  · intro v hv
    exact (Finset.mem_filter.mp hv).2
  · intro v₁ _ v₂ _ h
    simp only [Prod.mk.injEq] at h
    ext <;> omega
  · intro x hx
    refine ⟨(x.1 + a, x.2 + b), Finset.mem_filter.mpr ⟨hS x hx, ?_⟩, ?_⟩
    · simpa using hx
    · simp

theorem sum_shift (V S : Finset (ℤ × ℤ)) (a b : ℤ)
    (hS : ∀ x ∈ S, (x.1 + a, x.2 + b) ∈ V) :
    (∑ v ∈ V, if (v.1 - a, v.2 - b) ∈ S then (1 : ℤ) else 0) = S.card := by
  rw [Finset.sum_boole, card_shift V S a b hS]

theorem main : ViridisStatements.P1.BitQuad.Statement := by
  intro B
  set V := vertices B with hV
  -- inclusion facts
  have hB : ∀ x ∈ B, ∀ a b : ℤ, (a = 0 ∨ a = 1) → (b = 0 ∨ b = 1) →
      (x.1 + a, x.2 + b) ∈ V := by
    intro x hx a b ha hb
    rw [hV, mem_vertices]
    refine ⟨x, hx, ?_⟩
    rw [mem_corners]
    rcases ha with rfl | rfl <;> rcases hb with rfl | rfl <;> simp
  have hH : ∀ x ∈ hEdges B, ∀ a : ℤ, (a = 0 ∨ a = 1) → (x.1 + a, x.2 + 0) ∈ V := by
    intro x hx a ha
    rw [mem_hEdges] at hx
    rcases hx with hx | hx
    · have := hB _ hx a 0 ha (Or.inl rfl)
      simpa using this
    · have := hB _ hx a 1 ha (Or.inr rfl)
      simpa using this
  have hVe : ∀ x ∈ vEdges B, ∀ b : ℤ, (b = 0 ∨ b = 1) → (x.1 + 0, x.2 + b) ∈ V := by
    intro x hx b hb
    rw [mem_vEdges] at hx
    rcases hx with hx | hx
    · have := hB _ hx 0 b (Or.inl rfl) hb
      simpa using this
    · have := hB _ hx 1 b (Or.inr rfl) hb
      simpa using this
  -- the four expressions as sums over V
  have eV : ((V.card : ℤ)) = ∑ v ∈ V, (1 : ℤ) := by simp
  have eF11 := sum_shift V B 1 1 (fun x hx => hB x hx 1 1 (Or.inr rfl) (Or.inr rfl))
  have eF01 := sum_shift V B 0 1 (fun x hx => hB x hx 0 1 (Or.inl rfl) (Or.inr rfl))
  have eF10 := sum_shift V B 1 0 (fun x hx => hB x hx 1 0 (Or.inr rfl) (Or.inl rfl))
  have eF00 := sum_shift V B 0 0 (fun x hx => hB x hx 0 0 (Or.inl rfl) (Or.inl rfl))
  have eH10 := sum_shift V (hEdges B) 1 0 (fun x hx => hH x hx 1 (Or.inr rfl))
  have eH00 := sum_shift V (hEdges B) 0 0 (fun x hx => hH x hx 0 (Or.inl rfl))
  have eV01 := sum_shift V (vEdges B) 0 1 (fun x hx => hVe x hx 1 (Or.inr rfl))
  have eV00 := sum_shift V (vEdges B) 0 0 (fun x hx => hVe x hx 0 (Or.inl rfl))
  have eQ1 : ((Q1 B : ℕ) : ℤ) = ∑ v ∈ V, if windowCount B v = 1 then (1 : ℤ) else 0 := by
    rw [Finset.sum_boole]; rfl
  have eQ3 : ((Q3 B : ℕ) : ℤ) = ∑ v ∈ V, if windowCount B v = 3 then (1 : ℤ) else 0 := by
    rw [Finset.sum_boole]; rfl
  have eQD : ((QD B : ℕ) : ℤ) = ∑ v ∈ V, if isDiagonal B v = true then (1 : ℤ) else 0 := by
    rw [Finset.sum_boole]; rfl
  have key : 4 * euler B = ∑ v ∈ V,
      ((4 : ℤ) - 2 * ((if (v.1 - 1, v.2 - 0) ∈ hEdges B then 1 else 0)
        + (if (v.1 - 0, v.2 - 0) ∈ hEdges B then 1 else 0)
        + (if (v.1 - 0, v.2 - 1) ∈ vEdges B then 1 else 0)
        + (if (v.1 - 0, v.2 - 0) ∈ vEdges B then 1 else 0))
        + ((if (v.1 - 1, v.2 - 1) ∈ B then 1 else 0)
        + (if (v.1 - 0, v.2 - 1) ∈ B then 1 else 0)
        + (if (v.1 - 1, v.2 - 0) ∈ B then 1 else 0)
        + (if (v.1 - 0, v.2 - 0) ∈ B then 1 else 0))) := by
    simp only [Finset.sum_add_distrib, Finset.sum_sub_distrib, ← Finset.mul_sum]
    rw [eF11, eF01, eF10, eF00, eH10, eH00, eV01, eV00]
    unfold euler
    rw [← hV]
    simp
    ring
  rw [key, eQ1, eQ3, eQD, ← Finset.sum_sub_distrib, Finset.mul_sum,
    ← Finset.sum_sub_distrib]
  apply Finset.sum_congr rfl
  intro v hv
  obtain ⟨x, y⟩ := v
  have hne : (x - 1, y - 1) ∈ B ∨ (x, y - 1) ∈ B ∨ (x - 1, y) ∈ B ∨ (x, y) ∈ B := by
    rw [hV, mem_vertices] at hv
    obtain ⟨p, hp, hc⟩ := hv
    rw [mem_corners] at hc
    rcases hc with h | h | h | h <;> simp only [Prod.mk.injEq] at h <;> obtain ⟨rfl, rfl⟩ := h
    · right; right; right; exact hp
    · right; right; left; simpa using hp
    · right; left; simpa using hp
    · left; simpa using hp
  unfold windowCount isDiagonal bits
  simp only [sub_zero, mem_hEdges, mem_vEdges]
  by_cases ha : (x - 1, y - 1) ∈ B <;> by_cases hb : (x, y - 1) ∈ B <;>
    by_cases hc : (x - 1, y) ∈ B <;> by_cases hd : (x, y) ∈ B <;>
    simp [ha, hb, hc, hd] at hne ⊢

end Candidate
\end{lstlisting}

\section{Results}

\begin{itemize}
\item We give a Lean proof of Theorem~\ref{thm:main}, the sealed statement
  \texttt{P1.BitQuad}, for all finite rasters.
\item As an independent computational check, which is evidence and not a
  proof, the script \texttt{code/bitquad\_check.py} recomputes both sides for
  all $2^{16}=65{,}536$ subsets of a $4\times4$ box. It finds no
  counterexample (\texttt{results/bitquad\_4x4.json}). Euler characteristics
  from $-1$ to $4$ occur in that box.
\end{itemize}

\section{How L1 serves the L2/L3 bridge}

Let $p\notin B$ and $B'=B\cup\{p\}$. The windows whose bits differ between $B$
and $B'$ are exactly the four windows centred on the corners of $p$. By
Theorem~\ref{thm:main} applied to $B$ and to $B'$, $4(\chi(B')-\chi(B))$ equals
the change in $Q_1-Q_3-2Q_D$ over those four windows. That change depends only on the
$3\times3$ neighbourhood of $p$. This is the ``local half'' of L2. The non-local
half, the change in $C_8$ and in the number of enclosed 4-components of the
complement, still needs a separate argument: a pixel joining two
8-components can close off a new bounded region. L1 contributes nothing there.
The negative controls listed in the project context (the solid $2\times2$ block,
whose 8-neighbour graph $K_4$ has circuit rank 3 yet encloses nothing, and the
four-pixel N/S/E/W cross enclosing its centre) show that L2 must be stated
in terms of $\chi(K(B))$ and not in terms of the graph circuit rank.

\section{Limitations}

\begin{itemize}
\item The theorem concerns the combinatorial counts $V-E+F$ as defined in the
  sealed module. The development does not show that this number is the Euler
  characteristic of the topological space $K(B)\subset\mathbb{R}^2$. That
  identification is standard but is not formalised here, and L2/L3 are
  designed not to need it.
\item Only the 8-connected weighting of the bit quads is treated. The
  4-connected variant, with $+2Q_D$, is not formalised.
\item L2 and L3 remain open, so the \texttt{FaceCount} hypothesis of
  \emph{Trees Cannot Cut} is not yet discharged.
\end{itemize}

\section{Reproducibility}

The proof file is \texttt{proofs/P1.BitQuad.lean}. It imports only
\texttt{Mathlib} and the sealed module \texttt{ViridisStatements.P1.BitQuad}
(sha256 \texttt{cd3cdeba48f6\ldots{}8e23}), and it builds with toolchain
\texttt{leanprover/lean4:v4.28.0} and Mathlib commit
\texttt{8f9d9cff6bd728b17a24e163c9402775d9e6a365}. The check takes about 20
seconds. The brute-force script runs with the Python~3 standard library:
\texttt{python3 code/bitquad\_check.py results/bitquad\_4x4.json}. No external
data are used.

\section*{Changelog}
\begin{itemize}
\item 2026-10-28 (P1-T001): first version of the paper. Lean proof of L1
  (\texttt{P1.BitQuad}), proof architecture, brute-force check on the
  $4\times4$ box, and discussion of the role of L1 in L2/L3.
\end{itemize}

\bibliographystyle{plainnat}
\bibliography{refs}

\end{document}
